\section{Related Work}

The field of human motion generation has seen rapid advancements, driven by growing demand for scalable and diverse animation in virtual environments. Early approaches relied on manual animation or expensive motion-capture setups, which limited both the quantity and variety of available motion data~\cite{amass2019}. The emergence of data-driven methods—particularly deep generative models—has broadened possibilities for automated synthesis, but many frameworks remain constrained by the need for paired ground-truth data and are often limited to the action categories present in their training sets~\cite{edge2023, motiondiffuse2022}. Diffusion-based denoising models have recently improved fidelity and diversity~\cite{ddpm2020}, and music- or text-conditioned systems produce synchronized dances or language-aligned motions~\cite{motiondiffuse2022, edge2023, aist2021}; nevertheless, the dependence on curated 3D motion corpora sustains a persistent diversity bottleneck.

ViMo bypasses this bottleneck by framing video-to-motion as a generative task rather than strict reconstruction: a transformer-based diffusion model denoises 3D motion conditioned on 2D pose trajectories from casual video, comparing multi-view projections against the input to avoid explicit camera estimation and to tolerate occlusion and viewpoint shifts~\cite{vimo2024}. Building on this backbone, we focus on improving training dynamics: Min-SNR weighting accelerates convergence by reweighting loss terms according to signal-to-noise ratio, and vectorized timestep encoding replaces scalar diffusion indices with dense temporal embeddings to better capture rhythmic and long-range temporal structure~\cite{minsnr2024, ptss2024}. These modifications yield a model that trains efficiently on casual video and produces motions with sharper kinematic detail and improved rhythmic alignment.

\subsection{Action-Conditioned Generation}

Action-conditioned methods generate motions from categorical labels or seed poses, operating entirely within the training distribution. Degardin et al. propose a generative adversarial graph convolutional network that synthesizes human actions conditioned on discrete class labels, yet the model cannot extrapolate beyond the 10-15 action categories present during training \cite{gagcn2021}. Similarly, Action2Motion \cite{act2mo2020} employs a VAE-transformer architecture to generate 3D motions from action categories, but the generated repertoire remains bounded by the predefined set of 30 actions in the NTU RGB+D dataset. These approaches share a fundamental limitation: the conditioning signal is a one-hot or learned embedding of a seen category. The model learns a mapping from label to motion distribution but acquires no mechanism to compose unseen behaviors. The vocabulary of actions is closed at training time. Our video-to-motion framework sidesteps this restriction by conditioning on 2D pose trajectories rather than categorical tokens. A video of a never-before-seen dance style provides a kinematic blueprint that the diffusion model translates into plausible 3D motion, effectively enabling zero-shot generation of unseen categories.

\subsection{Pose Estimation from Videos}

Reconstructing 3D pose from monocular video follows a two-stage pipeline: detect 2D keypoints, then lift to 3D via convolutional or graph architectures. Cai et al. exploit spatial-temporal relationships through graph convolutional networks, aggregating joint correlations across frames to improve 3D accuracy \cite{gcnpose2019}. Martinez et al. demonstrate that a simple MLP operating on 2D detections can achieve competitive results when camera parameters are known \cite{poseest2017}. More recent work, SMPLer-X \cite{smpl2024}, scales training data to 4.5 million frames and predicts expressive SMPL-X parameters, yet relies on accurate camera calibration or approximates intrinsics from image metadata. These methods assume static cameras or easily estimable camera motion; under casual video with rapid zoom, pan, and shot cuts, the 2D-to-3D lifting objective becomes ill-posed. Frame-wise predictions accumulate drift, yielding jittery trajectories that violate physical plausibility. GLAMR \cite{glamr2022} attempts global trajectory optimization to smooth results, but computational cost is high and performance degrades under aggressive montage editing.

A related line of work synthesizes novel views of human interactions from sparse multi-view videos \cite{nvshismv2022}. This approach interpolates intermediate viewpoints given a few calibrated cameras, relaxing dense capture requirements. However, it still assumes known camera geometry and static scene capture, typically warping explicit 3D geometry or using light-field interpolation. Both strategies remain fragile under rapid camera motion or occlusions characteristic of casual footage. ViMo's generative formulation avoids explicit 3D regression entirely. Instead of minimizing per-joint reprojection error, the diffusion objective encourages synthesized motion to project consistently across multiple virtual views, implicitly handling camera variation without calibration.

\subsection{Temporal Modeling in Video Diffusion}
Standard diffusion models treat timestep as a scalar, ignoring frame-to-frame temporal structure and discarding periodicity and phase relationships inherent in motion~\cite{ddpm2020}. The Vectorized Timestep Approach redefines this by encoding timestep as a dense vector function of sequence position, allowing the transformer to capture long-range correlations and rhythmic patterns~\cite{ptss2024}. Our implementation introduces a \textbf{Probabilistic Timestep Sampling Strategy (PTSS)} to manage computational cost. In naive vectorized diffusion, sampling distinct timesteps per frame leads to a combinatorial explosion, with $1000^N$ combinations for $N$ frames. PTSS introduces a probability $p$ governing independent versus shared sampling: with probability $p$, each frame gets its own timestep for independent evolution; with probability $1-p$, a single timestep is broadcast across all frames, preserving efficiency. This hybrid balances temporal expressivity and training cost. In practice, $p=0.3$ balances diversity and stability, preventing overfitting to arbitrary frame-wise noise schedules while retaining fine-grained temporal control. Our experiments show PTSS improves beat-alignment scores for music-conditioned generation without increasing training time.

\subsection{Diffusion Training Optimization}
Training diffusion models requires careful balancing of loss contributions across the noise schedule. Ho et al. use constant weighting $w_t = 1$, treating all timesteps equally~\cite{ddpm2020}. SNR weighting $w_t = \alpha_t^2 / \sigma_t^2$ emphasizes high-SNR steps, but is numerically equivalent to constant weighting when predicting noise~\cite{impddpm2021}. Max-SNR-$\gamma$ clipping avoids zero weight at final timesteps but concentrates updates on low-noise regimes~\cite{pdfsdm2022}. Min-SNR-$\gamma$ weighting inverts this logic: $w_t = \min{\alpha_t^2 / \sigma_t^2, \gamma}$ caps the weight at small noise levels, forcing the model to prioritize high-noise timesteps where the signal structure is coarse~\cite{minsnr2024}. This accelerates early training and prevents overfitting to fine perturbations. Our implementation adopts $\gamma=5$ as default, following ablations in the original work. Empirically, Min-SNR reduces training epochs by 30\% and yields lower FID on kinetic features compared to constant weighting. For motion generation, where high-noise steps correspond to coarse pose and trajectory, this emphasis improves global structure before refining local joints. The weighting integrates seamlessly with the vectorized timestep encoder, forming a cohesive, efficient, and temporally aware training regime.